% ============================================================
% CAPITOLO 11 — Progetto 8: Setup Flutter e Prima App Mobile
% ============================================================
\chapter{Progetto 8: Setup Flutter e Prima App Mobile}\index{Flutter}\index{Dart}\index{Riverpod}\index{Dio}

\section*{Cosa Costruirai}
Un'applicazione mobile Flutter funzionante che:
\begin{itemize}
  \item Gira su Android e iOS da un unico codice
  \item Ha una schermata di login e una lista note (UI statica)
  \item Usa il pattern 0-code con \texttt{\_CONTEXT.md} per Flutter/Dart
\end{itemize}

\textbf{Tempo stimato}: 60--90 minuti\\
\textbf{Prerequisito}: Backend Notes API funzionante (Capitoli~6--10)

\begin{notebox}
\textbf{Box Tooling --- Stack scelto per questo esempio.}
\begin{itemize}
  \item \textbf{Framework:} Flutter~3.x (Dart~3.x)
  \item \textbf{State Management:} Riverpod~2.x
  \item \textbf{Routing:} GoRouter~14.x
  \item \textbf{HTTP Client:} Dio
\end{itemize}
\textbf{Alternative equivalenti:} React Native, Kotlin Multiplatform, .NET MAUI. Il \textbf{pattern} (UI dichiarativa, gestione stato, navigazione, chiamate API al backend) \`e comune a tutti i framework mobile cross-platform. Se scegli React Native, il metodo ADLC e il \texttt{\_CONTEXT.md} funzionano allo stesso modo.
\end{notebox}

% ---------------------------------------------------------
\section{Installare Flutter}

\begin{versionbox}
Le istruzioni di questo capitolo sono verificate con Flutter~3.x, Dart~3.x, Riverpod~2.x e GoRouter~14.x. L'ecosistema Flutter evolve rapidamente: i pacchetti cambiano API tra versioni minor. Se l'IA ti segnala un metodo deprecato o suggerisce una versione più recente di un pacchetto, fidati della correzione --- il pattern architetturale resta lo stesso. Esegui sempre \texttt{flutter pub get} dopo modifiche a \texttt{pubspec.yaml}.
\end{versionbox}

\begin{praticabox}[title={Pratica --- Installazione}]
\textbf{Windows:}
\begin{lstlisting}[language=bash]
winget install -e --id Google.Flutter
\end{lstlisting}

\textbf{macOS:}
\begin{lstlisting}[language=bash]
brew install --cask flutter
\end{lstlisting}

\textbf{Linux:}
\begin{lstlisting}[language=bash]
sudo snap install flutter --classic
\end{lstlisting}

\textbf{Verifica:}
\begin{lstlisting}[language=bash]
flutter doctor
\end{lstlisting}
Devi avere: Flutter SDK \checkmark, Android toolchain \checkmark, Connected device \checkmark.
\end{praticabox}

\begin{tipbox}
Per lo sviluppo Android senza Android Studio completo, puoi usare solo Android SDK command-line tools. Ma per i primi progetti, Android Studio con l'emulatore integrato è la scelta più semplice.
\end{tipbox}

Installa l'estensione \textbf{Flutter} (Dart-Code) per VS~Code --- abilita IntelliSense, esecuzione, debug, hot reload e widget inspector.

% ---------------------------------------------------------
\section{Creare il Progetto}

\begin{praticabox}[title={Pratica --- Scaffold del progetto}]
\begin{lstlisting}[language=bash]
cd notes-fullstack
flutter create notes_mobile
cd notes_mobile
\end{lstlisting}
Struttura generata:
\begin{lstlisting}[language={}]
notes_mobile/
+-- lib/
|   +-- main.dart          <- Punto di ingresso
+-- android/               <- Configurazione Android
+-- ios/                   <- Configurazione iOS
+-- test/                  <- Test
+-- pubspec.yaml           <- Dipendenze
+-- analysis_options.yaml  <- Regole di lint
\end{lstlisting}
\end{praticabox}

\begin{praticabox}[title={Pratica --- Verifica che funzioni}]
\begin{lstlisting}[language=bash]
flutter run
# Oppure per test rapidi:
flutter run -d chrome
\end{lstlisting}
\end{praticabox}

\begin{tipbox}
\texttt{flutter run -d chrome} è il modo più veloce per iterare. Non hai bisogno di un emulatore per la maggior parte del lavoro UI.
\end{tipbox}

% ---------------------------------------------------------
\section{Il Contesto per Flutter}

\begin{lstlisting}[language={},title={\texttt{\_CONTEXT.md} per Notes Mobile}]
# Progetto: Notes Mobile (Flutter)

## Scopo
App mobile Flutter collegata al backend Notes API.

## Tecnologie
- Flutter 3.x, Dart
- State Management: Riverpod 2
- HTTP: Dio
- Routing: GoRouter
- Storage locale: flutter_secure_storage
- UI: Material Design 3

## Architettura
lib/
  main.dart
  app.dart
  config/
    api_config.dart      <- URL backend, timeouts
    theme.dart           <- Tema Material Design 3
  models/
    user.dart / note.dart / api_response.dart
  services/
    api_service.dart     <- Client Dio configurato
    auth_service.dart    <- Login, logout, refresh
    notes_service.dart   <- CRUD note
  providers/
    auth_provider.dart   <- Stato autenticazione
    notes_provider.dart  <- Stato note
  screens/
    login_screen.dart
    dashboard_screen.dart
    note_detail_screen.dart
    note_form_screen.dart
  widgets/
    note_card.dart / empty_state.dart / error_banner.dart

## Convenzioni Dart/Flutter
- File: snake_case.dart
- Classi: PascalCase
- Variabili/funzioni: camelCase
- Widget: ogni widget in un file separato
- State: Riverpod (NON setState per stato condiviso)
- Modelli: classi immutabili con factory fromJson/toJson
- Max 3 livelli annidamento widget -> estrai

## Vincoli
- NON usare SharedPreferences per token JWT
- NON usare setState per stato condiviso tra schermate
- NON hardcodare URL. Usa api_config.dart
- Ogni schermata: loading, errore, vuoto, dati
- Material Design 3 con useMaterial3: true
\end{lstlisting}

% ---------------------------------------------------------
\section{\texttt{SKILL.md} per Flutter}

\begin{lstlisting}[language={},title={\texttt{SKILL.md} --- Flutter Mobile Developer}]
# Flutter Mobile Developer Skill

## Struttura Widget
class NoteCard extends ConsumerWidget {
  final Note note;
  final VoidCallback? onTap;
  const NoteCard({super.key, required this.note, this.onTap});

  @override
  Widget build(BuildContext context, WidgetRef ref) {
    return Card(
      child: ListTile(
        title: Text(note.title),
        subtitle: Text(note.content, maxLines: 2),
        onTap: onTap,
      ),
    );
  }
}

## Riverpod Pattern
final authProvider =
  StateNotifierProvider<AuthNotifier, AuthState>((ref) {
    return AuthNotifier(ref.read(authServiceProvider));
  });

## Design Material 3
- ColorScheme.fromSeed(seedColor: Colors.indigo)
- Theme.of(context).colorScheme per i colori
- Card con elevation: 0 e bordo per stile moderno
- FilledButton primario, OutlinedButton secondario
\end{lstlisting}

% ---------------------------------------------------------
\section{Generare la Struttura Base}

\begin{praticabox}[title={Pratica --- Genera l'app con Copilot}]
\begin{lstlisting}[language={}]
Leggi il _CONTEXT.md e il SKILL.md. Genera la struttura base.

Ordine:
 1. Dipendenze in pubspec.yaml: dio, riverpod,
    flutter_riverpod, go_router, flutter_secure_storage
 2. lib/config/api_config.dart
 3. lib/config/theme.dart (Material Design 3, seed: indigo)
 4. lib/models/ (Note, User, ApiResponse) immutabili
 5. lib/services/api_service.dart (Dio + interceptor)
 6. lib/providers/auth_provider.dart
 7. lib/app.dart con MaterialApp.router
 8. lib/screens/login_screen.dart (UI statica)
 9. lib/screens/dashboard_screen.dart (UI statica, dati mock)
10. GoRouter con redirect per protezione route
11. lib/main.dart
\end{lstlisting}
Dopo che l'IA ha aggiornato \texttt{pubspec.yaml}:
\begin{lstlisting}[language=bash]
flutter pub get
\end{lstlisting}
\end{praticabox}

% ---------------------------------------------------------
\section{Verifica e Iterazione}

\begin{praticabox}[title={Pratica --- Avvia l'app}]
\begin{lstlisting}[language=bash]
flutter run -d chrome
\end{lstlisting}
\end{praticabox}

\begin{center}
\begin{tabularx}{\textwidth}{l X}
\toprule
\textbf{Funzionalità} & \textbf{Test} \\
\midrule
App si avvia & Nessun errore di compilazione \\
Login screen & Bottoni Google/GitHub visibili \\
Tema & Colori Material~3 indigo \\
Router & Navigazione login $\leftrightarrow$ dashboard \\
Redirect & Accesso diretto a /dashboard $\rightarrow$ /login \\
\bottomrule
\end{tabularx}
\end{center}

\textbf{Hot Reload}: con l'app in esecuzione, modifica il testo di un bottone e salva. L'app si aggiorna \textbf{istantaneamente} --- il feedback loop più veloce dello sviluppo mobile.

\begin{checkpointbox}
App si avvia, mostra il login screen con tema Material~3, navigazione funziona.
\end{checkpointbox}

% ---------------------------------------------------------
\section{Il Mondo Dart in 10 Minuti}

Non hai bisogno di conoscere Dart per lavorare in 0-code. Ma riconoscere i pattern aiuta a revisionare:

\begin{lstlisting}[language=Java,title={Concetti chiave Dart}]
// Variabili: tipizzazione forte, null safety
String name = 'Notes App';
int? count;                       // Nullable
final items = <Note>[];           // Lista tipizzata

// Funzioni asincrone
Future<List<Note>> fetchNotes() async {
  final response = await dio.get('/api/notes');
  return (response.data['data'] as List)
      .map((json) => Note.fromJson(json))
      .toList();
}

// Classi
class Note {
  final String id;
  final String title;
  const Note({required this.id, required this.title});
}

// Widget = funzione che restituisce UI
class MyButton extends StatelessWidget {
  final String label;
  final VoidCallback onPressed;
  const MyButton({super.key,
    required this.label, required this.onPressed});

  @override
  Widget build(BuildContext context) {
    return FilledButton(
      onPressed: onPressed, child: Text(label));
  }
}
\end{lstlisting}

\begin{tipbox}
Dart assomiglia molto a Java/TypeScript. Se non li conosci, non preoccuparti: l'IA scrive il codice, tu verifichi che faccia quello che hai chiesto.
\end{tipbox}

% ---------------------------------------------------------
\section*{Riepilogo}

\begin{center}
\begin{tabularx}{\textwidth}{l X}
\toprule
\textbf{Aspetto} & \textbf{Dettaglio} \\
\midrule
Framework & Flutter~3.x con Dart \\
State & Riverpod~2 \\
Routing & GoRouter con redirect \\
Design & Material Design~3 \\
Struttura & Monorepo notes-fullstack/notes\_mobile/ \\
\bottomrule
\end{tabularx}
\end{center}
